\begin{helper}
Fix a history cell, a fixed blue support label \(B\), and a fixed red label
\(J\).  Let \(\mathcal T\) be the branch-label type, let \(\mathcal P\) be
the transcript family, and let
\[
\mathsf{Comp}(A,\beta,\tau)
\]
mean that the complete terminal assignment \(A\subseteq U\) is compatible
with the realized fixed-\((B,J)\) cell \((\beta,\tau)\).

Assume that compatibility is witnessed by an actual reconstructed sample
\(\omega_A\): for every compatible realized triple
\((A,\beta,\tau)\), the sample \(\omega_A\) belongs to the fixed-count target
and lies in the cell \(G_{B,J,\beta,\tau}\).  Assume further that the cell
and support-label construction identify the fixed labels with the present
staged-open supports of \(\omega_A\):
\[
B=
\{e\in\operatorname{StagedOpen}(\omega_A,\varepsilon,I):
  X_e(\omega_A)=1,\ e\text{ is blue}\},
\]
\[
R_{\beta,J}=
\{e\in\operatorname{StagedOpen}(\omega_A,\varepsilon,I):
  X_e(\omega_A)=1,\ e\text{ is red}\}.
\]
Finally assume that if \(e\in O_\beta\) is an open candidate and \(e\in A\),
then \(e\) is staged-open and present in the reconstructed sample
\(\omega_A\).

Then fixed-support exhaustion holds for compatible complete assignments: if
\(\mathsf{Comp}(A,\beta,\tau)\), \(e\in O_\beta\), and \(e\in A\), then
\[
e\in B\cup R_{\beta,J}.
\]
\end{helper}
\begin{proof}
Fix a compatible realized triple \((A,\beta,\tau)\) and an open candidate
\(e\in O_\beta\) with \(e\in A\).  By the reconstruction hypothesis there is
an associated target sample \(\omega_A\) in the fixed
\((B,J,\beta,\tau)\) cell.  The branch/transcript cell construction supplies
the two support identities displayed in the statement: \(B\) is the blue
present staged-open support of \(\omega_A\), and \(R_{\beta,J}\) is the red
present staged-open support of \(\omega_A\).  These identities are the
support-label information from
\(\noderef{FixedSetHistoryCellRedBranchTranscriptCells}\) and
\(\noderef{FixedSetHistoryCellRedSupportLabels}\), applied to the
reconstructed sample rather than only to an original target sample.

The open-candidate realization hypothesis says that \(e\) is staged-open and
present in \(\omega_A\).  There are two color cases.  If \(e\) is a blue
coordinate, then by the defining identity for \(B\), the staged-open and
present facts place \(e\) in \(B\).  If \(e\) is a red coordinate, then the
same facts and the defining identity for \(R_{\beta,J}\) place \(e\) in
\(R_{\beta,J}\).  In either case,
\[
e\in B\cup R_{\beta,J}.
\]

This proves the support-exhaustion input needed by
\(\noderef{FixedSetHistoryCellRedSelectedAbsenceSiblingPruning}\) from a
compatible assignment only after constructing the corresponding target cell
sample.  Thus the argument does not use support labels for unrelated
nonempty cells as if they applied to every terminal cylinder assignment.
\end{proof}
